\section{Verificación, validación y diseño estadístico}\label{sec:verificacion}

La Parte 1 fijó cinco criterios de verificación y validación: conservación,
orden temporal, capacidad de los recursos, casos límite y ritmo de paso por
torno. Como no hay datos reales de colas ni de salidas, la validación consiste en
comprobar que el modelo reproduce la comprobación previa de capacidad de la
Parte 1 y las distribuciones que esta fija, y que se comporta como predice la
teoría de colas. Antes se explica cómo se obtienen los indicadores y con qué
método estadístico se analizan.

\subsection{Indicadores}\label{sec:indicadores}

Los indicadores se calculan primero dentro de cada partido y después se analizan
entre partidos. Así no se trata a los miles de espectadores de una réplica, que
comparten colas y recursos, como si fueran observaciones independientes.

\begin{table}[htbp]
\centering\small
\caption{Cálculo de los indicadores dentro de una réplica.}
\label{tab:indicadores}
\begin{tabularx}{\textwidth}{@{}p{3.6cm}Yp{2.1cm}@{}}
\toprule
Indicador & Cálculo dentro de una réplica & Unidad \\
\midrule
Espera media y percentil 95 & Inicio del control menos llegada a la cola; media y percentil 95 de las esperas & min \\
Espera menor de 10~min & Porcentaje de personas con espera estrictamente menor de 10~min & \% \\
Cola máxima en una puerta & Mayor número de personas esperando a la vez en una puerta realmente usada & personas \\
Utilización de los tornos & Tiempo de servicio acumulado entre (tornos nominales $\times$ tiempo desde la apertura hasta el último control) & \% \\
Tiempo de acceso & Asiento menos llegada a la cola & min \\
Entrada tardía & De quienes llegan antes de $t=0$, porcentaje que se sienta después de $t=0$ & \% \\
Retraso por conflicto al resto del bloque & Diferencia del tiempo vomitorio--asiento entre otros espectadores expuestos y no expuestos a un conflicto en su bloque; excluye a quien lo causa & s por persona \\
Vaciado ordinario & Última salida menos final del partido & min \\
Máximo en una salida & Mayor número de personas a la vez en una salida, incluida la que la cruza & personas \\
Evacuación & Último en zona segura menos alarma; porcentaje a salvo a los 8~min & min; \% \\
\bottomrule
\end{tabularx}
\end{table}

Tres aclaraciones evitan lecturas erróneas:

\begin{itemize}
  \item La \textbf{entrada tardía} incluye control, marcha, vomitorio y posibles
    conflictos. La Parte 1 la llamaba «entrada tardía por colas», pero su
    definición no permite atribuir el retraso solo a la cola del torno; aquí se
    llama simplemente entrada tardía.
  \item La \textbf{utilización} promedia unas dos horas, desde la apertura hasta
    el último control. En X0 es del 40,43~\%, mientras que la intensidad de
    tráfico en la media hora punta es $\rho=0{,}93$: son ventanas y magnitudes
    distintas, y para la carga en la punta sirve $\rho$.
  \item El \textbf{percentil 95} publicado es la media entre réplicas de los
    percentiles de cada partido; la cola máxima es la media de los máximos por
    réplica. No se mezclan los espectadores de varias réplicas.
\end{itemize}

En P4 los mayores y las PMR se identifican con los mismos criterios en FIFO y
con prioridad, de modo que X0 y X3 comparan los mismos colectivos. En P7 se
distinguen la espera de todo el público y la del colectivo asignado a la puerta
A-1 que llega durante el cierre. En P8 el denominador del porcentaje a salvo es
la población que estaba dentro al sonar la alarma, aunque estuviera en tránsito;
quienes abandonan fuera o aún no han llegado no cuentan. Los indicadores que no
aplican a un escenario quedan sin valor y no cuentan en $n$.

\subsection{Método estadístico y campaña}\label{sec:campana}

\paragraph{Intervalos de confianza.} Para un indicador con valores
$Y_1,\dots,Y_n$ por réplica se calcula la media $\bar Y$ y el intervalo
t de Student del 95~\%:
\[
  \bar Y \pm h, \qquad h = t_{0{,}975;\,n-1}\,\frac{s}{\sqrt n}, \qquad
  s^2 = \frac{1}{n-1}\sum_{r=1}^n (Y_r-\bar Y)^2 .
\]
La semianchura $h$ describe la incertidumbre de la media entre réplicas, no la
dispersión entre espectadores ni un límite que todos los partidos respeten.

\paragraph{Diferencias emparejadas.} Para cada semilla común se calcula
$D_r = Y_{\text{alternativa},r} - Y_{\text{referencia},r}$ y se aplica la misma
fórmula a los $D_r$; $n$ es el número de parejas, que puede ser menor que el
tamaño de una de las muestras. Restar las medias de dos muestras completas no da
el mismo estimador. Las referencias son X0 para X1--X4 y S1, 6 acomodadores para
X5, sin redirección para X6 y todas las salidas para X7. Una diferencia negativa
reduce el indicador. Los IC son individuales para cada comparación; no son una
garantía simultánea para todos los indicadores.

\paragraph{Tamaño de muestra.} Se partió de 10 réplicas piloto por escenario y
se estimó el tamaño necesario para que la semianchura de la espera media no
superase el 5~\% de su valor:
$n = \lceil (t\,s/h_{\text{obj}})^2\rceil$ con $h_{\text{obj}} = 0{,}05\,|\bar Y|$,
con un mínimo de 30. Al terminar se recalculó el intervalo y se amplió la
muestra donde hacía falta. La precisión se fijó para la espera media de entrada,
no para cada indicador.

\begin{table}[htbp]
\centering\small
\caption{Réplicas definitivas por escenario y precisión alcanzada en la espera media.}
\label{tab:replicas}
\begin{tabular}{@{}lrr@{\hspace{2.5em}}lrr@{}}
\toprule
Escenario & $n$ & $h$ relativa (\%) & Escenario & $n$ & $h$ relativa (\%) \\
\midrule
X0 (base) & 180 & 4,34 & X6 (sin redirección) & 45 & 4,37 \\
X1 (2 tornos) & 30 & 1,93 & X6 (con redirección) & 199 & 3,71 \\
X1 (4 tornos) & 183 & 4,69 & X7 (todas las salidas) & 180 & 4,34 \\
X2 (llegada anticipada) & 138 & 4,97 & X7 (una salida bloqueada) & 180 & 4,34 \\
X3 (prioridad) & 181 & 4,32 & S1 (aforo $-$10~\%) & 247 & 4,19 \\
X4 (salida escalonada) & 180 & 4,34 & S1 (aforo $+$10~\%) & 73 & 4,50 \\
X5 (3 acomodadores) & 180 & 4,34 & S1 (velocidad $-$15~\%) & 180 & 4,34 \\
X5 (6 acomodadores) & 180 & 4,34 & S1 (velocidad $+$15~\%) & 180 & 4,34 \\
X5 (9 acomodadores) & 180 & 4,34 & S1 (servicio $-$25~\%) & 219 & 4,99 \\
 & & & S1 (servicio $+$25~\%) & 30 & 3,56 \\
\bottomrule
\end{tabular}
\end{table}

La campaña suma 2.965 réplicas en 19 escenarios: 13 configuraciones principales y
6 de sensibilidad. Todas cumplen el criterio, con una semianchura relativa máxima
del 4,9865~\%. Las 10 réplicas piloto de cada escenario forman parte del total,
sin contarse dos veces. Los escenarios que no cambian la entrada repiten la
espera de X0 semilla a semilla (apartado~\ref{sec:reproducibilidad}) y por eso
comparten su tamaño y precisión. El antiguo piloto de X6 con el cierre en
$t=+20$ es histórico y no se mezcla con la campaña.

\subsection{Conservación, orden temporal y capacidad}

\begin{itemize}
  \item \textbf{Conservación.} Las pruebas automáticas comprueban que cada
    espectador está en un único estado en cada transición, que es la igualdad de
    la Parte 1 persona a persona y no solo en la suma. En los datos definitivos,
    en las 2.605 réplicas sin alarma todos los espectadores generados entran, se
    sientan y salen, sin ninguno pendiente; en las 360 réplicas de X7 todos quedan
    a salvo (a $t=+40$ ya no hay cola exterior y nadie abandona).
  \item \textbf{Orden temporal.} Se comprueba llegada $\le$ inicio del control
    $\le$ fin del control $\le$ llegada al vomitorio $\le$ asiento $\le$ salida.
  \item \textbf{Capacidad.} Ningún torno atiende a más personas que su
    capacidad, también con averías; nunca hay más de 8 personas en un vomitorio,
    un conflicto retiene su plaza, cada macrozona usa solo sus acomodadores y cada
    salida deja pasar una persona cada 0,227~s.
\end{itemize}

Las tablas agregadas de la campaña no permiten reconstruir cada transición de las
2.965 réplicas; esa evidencia procede de la auditoría hecha al simular, de las
pruebas automáticas y de las reproducciones individuales.

\subsection{Capacidad de los tornos}

La Parte 1 estimó la capacidad real de un torno en 13,6~pers/min: servicio medio
de $4 + 0{,}03\times 10 = 4{,}3$~s con la relectura, y un 97,4~\% del tiempo
operativo por las averías ($240/(240+6{,}5)$). El cálculo analítico reproduce
esas cifras: servicio medio 4,30~s, disponibilidad 0,974, 13,59~pers/min por
torno y 978~pers/min con los 72 tornos. En los datos de X0 ($n=180$) se
comprueban los dos componentes por separado:

\begin{table}[htbp]
\centering\small
\caption{Componentes de la capacidad en X0 ($n=180$).}
\label{tab:componentes}
\begin{tabular}{@{}lrr@{}}
\toprule
Componente & Teórico & Modelo (IC 95~\%) \\
\midrule
Espectadores con relectura & 3,00~\% & 3,001 $\pm$ 0,011~\% \\
Averías por réplica (72 tornos, de la apertura al final) & 59,4 & 60,1 $\pm$ 1,0 \\
\bottomrule
\end{tabular}
\end{table}

Ambos valores teóricos quedan dentro del intervalo. El ritmo resultante se midió
en los periodos con cola sostenida, en los que los tornos no paran por falta de
público: en una réplica de X1 con 2 tornos, idéntica a la de la campaña, salen
13,54~pers/min por torno, un 0,4~\% por debajo de los 13,59 teóricos. Sin
incidencias, el ritmo medido es de 14,98~pers/min (teórico 15,00). Los
vomitorios admiten 411~pers/min por macrozona frente a los 245 que dejan pasar los
tornos, de modo que no son un cuello de botella en la entrada sin conflictos.

\subsection{Intensidad de tráfico y espera}\label{sec:rho}

La intensidad de tráfico es $\rho = \lambda_{\text{pico}}/\text{capacidad}$, con
$\lambda_{\text{pico}}$ la tasa de llegadas en la franja más cargada. Para
contrastarla con la espera observada se construyó un \emph{modelo fluido} de cada
puerta: llegadas y servicio como caudales continuos, con la capacidad media de
13,59~pers/min por torno. Ignora la variabilidad: con $\rho<1$ predice espera
nula y con $\rho>1$ una espera que crece a ritmo $\rho-1$ (minutos de espera por
minuto de reloj) mientras dura la sobrecarga y que después se reduce hasta
vaciarse. La simulación debe quedar algo por encima.

\begin{table}[htbp]
\centering\small
\caption{Intensidad de tráfico frente a la espera (min). Espera fluida: modelo
determinista; observada: media de las réplicas con IC del 95~\%.}
\label{tab:rho}
\begin{tabular}{@{}lrrrrr@{}}
\toprule
Escenario & $n$ & $\rho$ & Espera fluida & Espera observada & Percentil 95 observado \\
\midrule
X1 (4 tornos) & 183 & 0,70 & 0 & 0,016 $\pm$ 0,001 & 0,076 $\pm$ 0,002 \\
X2 (llegada anticipada) & 138 & 0,81 & 0 & 0,050 $\pm$ 0,003 & 0,209 $\pm$ 0,012 \\
S1 (aforo $-$10~\%) & 247 & 0,84 & 0 & 0,090 $\pm$ 0,004 & 0,489 $\pm$ 0,032 \\
X0 (base) & 180 & 0,93 & 0 & 0,192 $\pm$ 0,008 & 1,150 $\pm$ 0,068 \\
S1 (aforo $+$10~\%) & 73 & 1,02 & 0,205 & 0,486 $\pm$ 0,022 & 2,438 $\pm$ 0,130 \\
S1 (servicio $+$25~\%) & 30 & 1,15 & 1,304 & 1,402 $\pm$ 0,050 & 4,919 $\pm$ 0,211 \\
X1 (2 tornos) & 30 & 1,40 & 3,826 & 4,016 $\pm$ 0,078 & 11,764 $\pm$ 0,266 \\
\bottomrule
\end{tabular}
\end{table}

\begin{itemize}
  \item \textbf{La espera crece con $\rho$.} En los ocho escenarios sin cierre y
    con entrada distinta, el orden de las esperas medias es exactamente el de
    $\rho$.
  \item \textbf{Con $\rho$ muy por debajo de 1} (X1 con 4 tornos, X2) la espera
    media es de 1 a 3~s: lo observado es solo la variabilidad.
  \item \textbf{Con $\rho \approx 0{,}93$} (X0) la espera media es de 11,5~s y el
    percentil 95 de 1,15~min. Hay picos de cola por la variabilidad, pero no
    acumulación sostenida.
  \item \textbf{Con $\rho \approx 1$} (aforo $+$10~\%) el fluido predice
    0,21~min y se observan 0,49: en el entorno de $\rho=1$ domina la variabilidad
    y el fluido es solo una cota inferior.
  \item \textbf{Caso límite saturado} (X1 con 2 tornos, $\rho=1{,}40$). La espera
    crece sin cota durante la punta ($-30$ a 0~min) y se vacía después. El fluido
    predice una espera máxima de $(\rho-1)\times 30 \approx 11{,}9$~min, dentro del
    intervalo del percentil 95 observado (11,764~$\pm$~0,266~min). La espera media
    coincide con una desviación del 5~\% (3,826 fluido frente a
    4,016~$\pm$~0,078 observado), y con el servicio $+$25~\% la desviación es
    del 7~\%; en ambos casos lo observado queda por encima, como se espera.
    Que un modelo construido con 13,59~pers/min reproduzca la espera de los
    escenarios saturados confirma de forma independiente la capacidad efectiva.
  \item \textbf{$\rho$ de X2 = 0,81.} La espera observada en X2 (0,050~min) queda
    entre la de $\rho=0{,}70$ (0,016) y la de $\rho=0{,}84$ (0,090), coherente con
    0,81 y no con el 0,65 de la Parte 1.
  \item \textbf{Cierre de puerta en $t=-30$} (X6). Para quienes llegan a A-1
    durante el cierre, el fluido predice 9,66~min sin redirección y se observan
    9,626~$\pm$~0,120 ($n=45$), dentro del intervalo. Con redirección predice
    0,59~min y se observan 0,543~$\pm$~0,026 ($n=199$): coincide con una
    desviación del 8~\%, en este caso por debajo, porque ir a la puerta con menos
    cola reparte mejor que el reparto igual que supone el fluido.
\end{itemize}

\subsection{Distribuciones de entrada}

Se regeneraron, sin simular, los espectadores de las réplicas de X0 (180) y X2
(138) con sus semillas, y se contrastaron con las proporciones de la Parte 1
mediante $\chi^2$:

\begin{table}[htbp]
\centering\small
\caption{Contraste $\chi^2$ de los repartos de entrada con la Parte 1.}
\label{tab:repartos}
\begin{tabular}{@{}lrrrr@{}}
\toprule
Reparto & Réplicas & \makecell[r]{Desviación\\máxima} & \makecell[r]{$p$ ($\chi^2$\\agregado)} & \makecell[r]{Réplicas\\con $p<0{,}05$} \\
\midrule
Perfil por franjas, X0 & 180 & 0,010 p.\,p. & 0,62 & 3,3~\% \\
Perfil por franjas, X2 & 138 & 0,018 p.\,p. & 0,37 & 3,6~\% \\
Grupos 88/10/2~\%, X0 & 180 & 0,015 p.\,p. & 0,28 & 5,0~\% \\
Macrozonas 25~\%, los 19 escenarios & 2.965 & 0,006 p.\,p. & 0,21 & 5,5~\% \\
\bottomrule
\end{tabular}
\end{table}

Ningún reparto se rechaza y la fracción de réplicas rechazadas al 5~\% ronda el
5~\%, lo esperado si las proporciones son las de la Parte 1. Las réplicas
regeneradas dan exactamente los recuentos por macrozona de la campaña.

\subsection{Salida y evacuación frente a la Parte 1}

Para aislar el paso por las salidas, una prueba usa una macrozona ($N/4$, unas
12.194 personas) que llega junta a las salidas, con 1~m de recorrido y 1~s de
reacción:

\begin{table}[htbp]
\centering\small
\caption{Paso por las salidas de una macrozona frente a la Parte 1.}
\label{tab:salidas}
\begin{tabular}{@{}llrrr@{}}
\toprule
Caso & Salidas & Paso por persona & Parte 1 & Modelo \\
\midrule
Salida ordinaria & 6 de 4~m & 0,227~s & 7,7~min & 7,7~min \\
Evacuación & 6 de 5,2~m & 0,175~s & 5,9~min & 5,9~min \\
Evacuación con una salida bloqueada & 5 de 5,2~m & 0,175~s & 7,1~min & 7,1~min \\
\bottomrule
\end{tabular}
\end{table}

Estas cifras miden solo el paso por las salidas. Los tiempos del partido completo
(sección~\ref{sec:resultados}) son mayores porque incluyen la marcha, la reacción
y el reparto aleatorio del público; no se han ajustado parámetros para que ambos
coincidan.

\subsection{Casos límite}

\begin{table}[htbp]
\centering\small
\caption{Casos límite comprobados con pruebas automáticas o réplicas específicas.}
\label{tab:limites}
\begin{tabularx}{\textwidth}{@{}p{5cm}Y@{}}
\toprule
Caso & Resultado \\
\midrule
Muchos tornos, sin incidencias & Espera prácticamente nula (percentil 95 $<$ 0,01~min) y exactamente 0 con la misma demanda \\
2 tornos ($\rho=1{,}40$) & La espera crece en la punta con pendiente $\rho-1\approx 0{,}40$~min por minuto (tolerancia 15~\%) \\
4 tornos ($\rho=0{,}70$) & Sin tendencia en la punta y espera media $<$ 0,1~min \\
2, 3 y 4 tornos & El percentil 95 y la utilización decrecen con el número de tornos \\
Cierre en $-30$ sin redirección & Espera de los afectados próxima al fluido ($\approx$ 9,7~min) \\
Cierre en $-30$ con redirección & Espera $<$ 1,5~min y al menos 5 veces menor que sin ella \\
Acomodadores casi ilimitados (5 réplicas) & La duración media de un conflicto baja a 4,62~min, la media teórica de la resolución (4,67): la saturación con pocos acomodadores se debe a que faltan, no a un error del modelo \\
X7 con todas las salidas, réplicas que no vacían en 8~min & Solo llega tarde una persona, siempre PMR y sin cola en la salida; sin las PMR, el vaciado nunca pasa de 7,36~min \\
\bottomrule
\end{tabularx}
\end{table}

\subsection{Reproducibilidad y números aleatorios comunes}\label{sec:reproducibilidad}

\begin{itemize}
  \item \textbf{Misma semilla, mismos resultados.} Se repitió una réplica de cada
    uno de los 19 escenarios con su semilla original: los 21 indicadores y los
    datos de auditoría coinciden en todos los casos. Repetir una réplica comprueba
    la reproducibilidad, pero no añade réplicas a la campaña.
  \item \textbf{Números comunes.} En los escenarios que no modifican la entrada
    (X4, X5 con 3, 6 y 9 acomodadores, las dos variantes de X7 y S1 de
    velocidad), los indicadores de entrada (espera media, percentil 95, cola
    máxima, relecturas) son idénticos a los de X0 en las 180 semillas comunes. El
    número de relecturas coincide también en X1, X2, X3 y X6.
  \item \textbf{Efecto sobre la precisión.} La correlación entre escenarios
    emparejados es siempre positiva, así que emparejar nunca empeora el intervalo
    de una diferencia. En X6 (con frente a sin redirección, correlación 0,94) la
    varianza de la diferencia baja al 7~\% de la que tendría con muestras
    independientes y el intervalo se estrecha 3,7 veces; entre X1 o X2 y X0 la
    ganancia es pequeña (de 1,07 a 1,20 veces), porque sus varianzas son muy
    distintas.
\end{itemize}

\subsection{Validación frente a la Parte 1}

\begin{table}[htbp]
\centering\small
\caption{Validación de coherencia con la comprobación previa de la Parte 1.}
\label{tab:validacion}
\begin{tabularx}{\textwidth}{@{}p{4.2cm}p{2.7cm}Yp{2.2cm}@{}}
\toprule
Magnitud & Parte 1 & Modelo & ¿Coincide? \\
\midrule
Servicio medio con relectura & 4,3~s & 4,30~s; relectura en el 3,001~$\pm$~0,011~\% & Sí \\
Torno operativo & 97~\% & 97,4~\%; 60,1~$\pm$~1,0 averías frente a 59,4 & Sí \\
Capacidad por torno & 13,6~pers/min & 13,59; 13,54 medido con cola & Sí \\
Capacidad total & 978~pers/min & 978 & Sí \\
$\rho$ con 2 / 3 / 4 tornos & 1,40 / 0,93 / 0,70 & 1,40 / 0,93 / 0,70 & Sí \\
$\rho$ de X2 & 0,65 & 0,81 & No: error de la Parte 1 \\
Vomitorios frente a tornos & 411 frente a 245~pers/min & 411 y 245 & Sí \\
Repartos de entrada & 25~\%; perfil; 88/10/2~\% & $\chi^2$, $p \ge 0{,}21$ & Sí \\
Paso por salidas & 7,7 / 5,9 / 7,1~min & 7,7 / 5,9 / 7,1~min & Sí \\
Instante del cierre & $t=+20$~min & $t=-30$~min & No: cambio decidido \\
Ritmo por torno según la fuente F2 & unos 4~s, 15~pers/min & 14,98~pers/min sin incidencias & Sí con la cifra que cita la Parte 1 \\
\bottomrule
\end{tabularx}
\end{table}

\paragraph{Límites de la validación.} La Parte 1 pedía comprobar que el ritmo por
torno coincide con la fuente F2. El contraste se ha hecho con la cifra que la
Parte 1 atribuye a F2 (unos 4~s por persona, 15~pers/min por torno); esa cifra no
se ha vuelto a comprobar en el artículo original, así que el contraste no es una
validación independiente de F2. En
general, la validación comprueba que el modelo hace lo que la Parte 1 especifica
y lo que predice la teoría de colas; no demuestra que reproduzca las colas reales
de San Mamés. El perfil de llegadas y los parámetros de las incidencias son
hipótesis: la verificación muestra que el modelo los aplica bien, no que sean
ciertos.
